\newpage{\ }
\cleardoublepage

\chapter[Capítulo 5. Discusión general]{Discusión general}\label{cap:discu}

En esta tesis avanzamos el entendimiento de los factores que controlan el fuego en el norte de la Patagonia Andina argentina a diversas escalas temporales y espaciales, aportando evidencia a partir de un estudio de campo y del análisis de la ocurrencia de incendios a escala de paisaje. En particular, pudimos separar efectos del tipo de vegetación de factores ambientales asociados. Finalmente, integramos lo aprendido en modelos de simulación de incendios, que permitieron proyectar la incidencia del fuego en la región bajo escenarios climáticos futuros. 

Un factor común a lo largo de esta tesis fue la contrastante inflamabilidad entre bosques y matorrales, observada a distintas escalas. A escala de sitio, los matorrales presentaron mayor probabilidad de ignición de combustibles muertos debido a su menor contenido de humedad comparado con los bosques (Fig.~\ref{fig:cap2-06}); a escala de paisaje, los matorrales mostraron mayor probabilidad de quema comparados con comunidades boscosas, incluso controlando por factores ambientales correlacionados (Fig.~\ref{fig:cap3-06}C). Estas diferencias también emergieron al modelar por separado la ignición, escape y propagación del fuego: los matorrales presentaron mayor probabilidad de ignición y de propagación que los bosques (Fig.~\ref{fig:ig_spat_veg} y ~\ref{fig:spread-veg-effects}). Estudios previos ya habían reportado diferencias estructurales que harían a los matorrales más inflamables que los bosques \citep{paritsis_positive_2015, blackhall_effects_2017, tiribelli_changes_2018}, y a escala de paisaje, los matorrales mostraban mayor probabilidad de quema \citep{mermoz_landscape_2005, paritsis_habitat_2013, morales_stochastic_2015}. Sin embargo, estos estudios presentaban algunas limitaciones para evaluar diferencias en inflamabilidad entre estos tipos de vegetación. Por un lado, las diferencias estructurales entre matorrales y bosques podrían perder relevancia bajo condiciones atmosféricas extremas \citep{Podur2009, Parks2015, Collins2019, jones_extreme_2023}, y por otro lado, la mayor probabilidad de quema de los matorrales observada a escala de paisaje podría estar explicada, al menos en parte, por otros factores correlacionados con el tipo de vegetación, como altitud -que influye en la temperatura-, precipitación, y uso humano. En esta tesis evaluamos estos aspectos, y encontramos que (1) los tipos de vegetación presentan inflamabilidad contrastante incluso bajo condiciones atmosféricas extremas (Fig.~\ref{fig:ig_spat_veg}), y (2) aún controlando por factores físicos o de uso humano a escala de paisaje, los matorrales sobresalen como una de las comunidades más inflamables (Fig.~\ref{fig:cap3-06}C). A su vez, encontramos que los bosques pueden presentar inflamabilidad similar a la de los matorrales cuando tienen baja productividad (bajo NDVI; Fig.~\ref{fig:cap3-05}B y \ref{fig:FI}A), probablemente porque presentan doseles abiertos, lo que les confiere características similares a las de los matorrales en cuanto a humedad y continuidad del combustible. 

La mayor inflamabilidad de los matorrales y bosques de baja productividad tiene implicancias relevantes para el fuego si consideramos los efectos del cambio climático sobre la vegetación. Ante un clima más cálido y seco, la regeneración posfuego de los bosques se verá limitada \citep{tercero-bucardo_field_2007}, aumentando la probabilidad de que sean reemplazados por matorrales a altitud baja e intermedia y por pastizales en la franja altitudinal superior \citep{kitzberger_firevegetation_2016, kitzberger_projections_2022}. A su vez, estas condiciones climáticas promueven el decaimiento y la mortalidad de árboles (ya sea parcial o total; \citealt{suarez2004factors,gutierrez2014increased,rodriguez2016influence,molowny2017changes}), con lo cual, los bosques pueden volverse más abiertos. Esto aumentaría la tasa de desecación de combustibles muertos y la continuidad de combustibles finos en el sotobosque, aumentando la inflamabilidad. Además, aumentaría la cantidad de combustible muerto en pie. De esta manera, sería esperable que el paisaje se vuelva más inflamable. Sin embargo, hay que tener en cuenta que con el cambio climático el aumento de $\text{CO}_2$ atmosférica podría aumentar la eficiencia hídrica de las plantas, pudiendo atenuar, los efectos del clima sobre la mortalidad \citep{silva2013probing,sperlich2020gains}.

El notable efecto de la vegetación sobre el fuego sugiere que, en las zonas accesibles, se podrían aplicar estrategias de manejo para reducir la inflamabilidad del paisaje. En particular, promover la formación de doseles altos y cerrados podría contribuir a mantener una mayor humedad del combustible  y reducir su continuidad en el sotobosque, disminuyendo así la probabilidad de ignición y propagación del fuego. Esto podría lograrse a través del manejo de matorrales, promoviendo el desarrollo de árboles de un solo fuste, y plantando individuos de especies de gran porte, como coihue, lenga y ciprés, u otras de menor porte pero con elevado contenido de humedad, como maitén. A su vez, la reducción de la biomasa del sotobosque mediante la extracción controlada de leña y madera, en combinación con el pastoreo de ganado, podría contribuir a disminuir la carga de combustible fino. Sin embargo, se deben seguir desarrollando e investigando estrategias de manejo para compatibilizar objetivos de disminución del peligro de incendios y rentabilidad económica \citep{peri_review_2016, goldenberg_efecto_2018, goldenberg_effects_2020, goldenberg_shrubland_2021}. Por otro lado, la elevada inflamabilidad de las plantaciones de coníferas exóticas y de los sitios que estas invaden resalta la necesidad de manejar los combustibles de las plantaciones y controlar las invasiones para disminuir la inflamabilidad del paisaje \citep{paritsis_pine_2018, cingolani2023enrichment, franzese_legacy_2022}.

Más allá del notable control del fuego por parte de la vegetación, las condiciones atmosféricas resultaron ser uno de los factores más influyente en la ocurrencia de incendios (Fig.~\ref{fig:cap3-03}, \ref{fig:igrate}, \ref{fig:spread-params}, y \ref{fig:proj-means}). Con el cambio climático, las condiciones atmosféricas propicias para el fuego serán cada vez más frecuentes y severas. A su vez, esto se agravará con la creciente incidencia de rayos, los cuales generan focos en zonas remotas que resultan muy difíciles de controlar (\citeauthor{kitzbergerLightning}, en revisión). Ante este escenario, será necesario invertir en estrategias de detección y supresión temprana de focos generados por rayos antes de que las condiciones atmosféricas generen incendios de gran tamaño e intensidad que limiten su control. A su vez, reducir considerablemente las igniciones humanas sería altamente beneficioso, para que los recursos destinados a la supresión del fuego puedan destinarse principalmente a los inevitables incendios generados por rayos. Estudios en otras regiones sugieren que las igniciones por humanos pueden ser reducidas notablemente mediante campañas educativas \citep{hesseln_wildland_2018}. Sin embargo, algunos incendios intencionales ocurren bajo condiciones atmosféricas extremas buscadas por sus perpetradores con el fin de generar un gran impacto. Esto sugiere que las tensiones entre distintos actores sociales deberían ser consideradas a la hora de buscar reducir las igniciones por humanos \citep{buhler2013demography}.

Además del fuerte control climático sobre la incidencia del fuego, en el Capítulo~\ref{cap:patrones} encontramos una marcada variabilidad espacial espacial en la incidencia del fuego (Tabla~\ref{tab:cap3-spatvar}). Estas diferencias tan grandes en la proporción del paisaje quemada en las distintas regiones probablemente no se explique únicamente por sus diferencias en topografía y vegetación, sino que podrían estar asociadas a diferencias en el manejo del fuego. De ser así, evaluar en detalle por qué algunas regiones se queman más que otras, en relación a los recursos destinados a la prevención y control del fuego, y al uso intencional del fuego por parte de las personas podría informar estrategias relevantes para disminuir la incidencia del fuego. 

Pero incluso considerando mejores estrategias para prevenir y controlar el fuego, el aumento previsto en la incidencia del fuego a causa del cambio climático y la dificultad para disminuir las igniciones humanas sugiere que estas estrategias probablemente no sean suficientes. Además, se requerirán estrategias para atenuar el inevitable impacto del fuego sobre las sociedades. En particular, esto podría lograrse implementando estrategias de evacuación ante incendios \citep{cova2002microsimulation}, promoviendo la construcción de infraestructura poco inflamable \citep{quarles2010home} y manejando la vegetación en las cercanías de la infraestructura humana, disminuyendo la biomasa de combustibles y plantando especies con un alto contenido de humedad \citep{reinhardt2008objectives, syphard2014role, blackhall_flammability_2019}. A su vez, a una escala temporal mayor, se podría diseñar un paisaje en donde las viviendas se acumulen en zonas poco vegetadas en vez de seguir desarrollando una amplia interfaz urbano-forestal en donde las viviendas se encuentran altamente expuestas al fuego \citep{syphard2013land, moritz2014learning, barrett2019reducing}. Sin embargo, políticas de este tipo probablemente sean más difíciles de implementar, requiriendo una fuerte presencia del Estado.   

El estudio del fuego y la vegetación con un enfoque de modelado espacialmente explícito puede realizar importantes contribuciones al manejo de los ecosistemas del norte de la Patagonia Andina argentina. En particular, las proyecciones de la dinámica de la vegetación y el fuego en esta región probablemente sean más atinadas avanzando nuestro entendimiento de los siguientes aspectos: 

\begin{itemize}
    \item factores que controlan la incidencia de rayos y cómo la misma se verá afectada por el cambio climático;
    \item efectos de la sequía sobre la inflamabilidad de los bosques;
    \item dinámica de la vegetación (incluyendo regeneración posfuego) ante un clima más cálido y seco;
    \item dinámica de floración de la caña \textit{Chusquea culeou}, cuyos eventos de floración y seguida mortalidad sincronizada aumentan la inflamabilidad del paisaje dramáticamente;
    \item dinámica de invasión de coníferas exóticas y sus impactos en la inflamabilidad a escala de paisaje;
    \item herbivoría por grandes mamíferos (ganado y especies invasoras), que puede disminuir la carga combustible pero también limitar la regeneración posfuego de los bosques;
    \item factores que regulan la marcada variabilidad espacial en la incidencia del fuego, probablemente asociados a aspectos del uso humano del territorio, y
    \item factores que regulan la tasa de ignición humana, y la efectividad de políticas para reducirlas.
\end{itemize}

\noindent Integrando estos conocimientos en modelos espacialmente explícitos podremos explorar estrategias de manejo que contribuyan a generar paisajes más resilientes al fuego. Los modelos de simulación de incendios elaborados en esta tesis presentan un buen punto de partida, habiéndose desarrollado en un lenguaje de código abierto que facilita su continuo desarrollo.